\section{Related Work}

\textbf{Learning to land.} Deep RL has landed quadrotors on ground markers, with one deep Q-network per stage and domain randomization \cite{polvara2017,polvara2020}, and on moving platforms \cite{rodriguez2018}, also in wind \cite{peter2024}. WaveLander \cite{li2026wavelander} is hierarchical: its policy outputs only a vertical velocity reference, and a conventional controller holds attitude and lateral position. Amendola et al. \shortcite{amendola2024} review the field. These studies learn the approach and touchdown. In our stack a conventional controller does that, and learning decides when and where to descend around other agents.

\textbf{Safety layers and perception.} NavRL \cite{xu2025navrl} trains a proximal policy optimization (PPO) policy \cite{schulman2017} among static and moving obstacles, behind a velocity-obstacle safety shield. Shielding in general blocks a learned policy's unsafe actions with a separate check \cite{alshiekh2018}. VO-Safe RL \cite{lin2024vosafe} shares our structure most closely. Its PPO policy picks waypoints for a low-level controller and is penalized when visual odometry fails, but its scenes have no moving obstacles. Perception-aware learners fly agile paths in clutter \cite{song2023} or slow down where predictions are uncertain \cite{kahn2017}. Planners reduce state-estimation uncertainty \cite{zhang2018} or keep a target in view \cite{falanga2018}, and policies can adapt to the localization precision each region needs \cite{puthumanaillam2024}. Our GNSS-denied flights use absolute visual localization, reviewed by Couturier and Akhloufi \shortcite{couturier2024}.

\textbf{Wind.} GustPilot \cite{habel2026gustpilot} pairs a learned velocity-reference policy with an incremental nonlinear dynamic inversion controller that rejects gusts. Residual RL adds a learned correction to a conventional controller \cite{johannink2019,silver2018}. Our flying learner is residual in that sense.

\textbf{Where our results sit.} Learned wind control has been shown up to 12.1~m/s on hardware for trajectory tracking \cite{oconnell2022} and up to 15~m/s in simulation \cite{altasim2026}. We test landing up to 20~m/s, in simulation only. Without learning, Falanga et al. \shortcite{falanga2020} dodged thrown objects, which do not steer, more than 90\% of the time outdoors. Apart from field reports of magpies chasing drones \cite{lyons2018}, we found no study of bird attacks on delivery drones. Benali and Hamoud \shortcite{benali2026} report 52 to 92\% avoidance of a helicopter on one fixed path, without a planner baseline. Elsewhere, a planner or a fixed search pattern beat the learned policy in static scenes \cite{xu2025navrl,yao2026}, as our ladder predicts for unmapped obstacles.
